\documentclass[12pt]{article}

\usepackage[a4paper, top=2.5cm, bottom=2.5cm, left=3cm, right=3cm]{geometry}
\usepackage[brazil]{babel}
\usepackage[utf8]{inputenc}
\usepackage[T1]{fontenc}
\usepackage{graphicx}
\usepackage{url}
\usepackage{booktabs}
\usepackage{array}
\usepackage{tabularx}
\usepackage{multirow}
\usepackage{listings}
\usepackage{xcolor}
\usepackage{hyperref}

% Define resumo environment (substitui sbc-template.sty)
\newenvironment{resumo}{%
  \vspace{1em}%
  {\centering\textbf{Resumo}\par}%
  \vspace{0.5em}%
  \small%
}{%
  \par\normalsize\vspace{1em}%
}

% -------------------------------------------------------
% Estilo para blocos de código / arquitetura
% -------------------------------------------------------
\lstset{
  basicstyle=\ttfamily\small,
  breaklines=true,
  frame=single,
  backgroundcolor=\color{gray!10}
}

\sloppy

% -------------------------------------------------------
% TÍTULO, AUTORES, INSTITUIÇÃO
% -------------------------------------------------------
\title{AgroLeilões: Desenvolvimento e Avaliação de uma Plataforma Web\\
       para Leilões Agropecuários em Tempo Real}

\author{Matheus [SOBRENOME]%
  \thanks{Engenharia de Software, ICEV -- Instituto de Ensino Superior do Ceará.
  Projeto de Extensão VII -- Pesquisa e Inovação (2026.1).
  Orientador: Prof. Luciani Vieira.
  \texttt{[SEU-EMAIL]@somosicev.com}}}

% -------------------------------------------------------
\begin{document}

\maketitle

% -------------------------------------------------------
% RESUMO
% -------------------------------------------------------
\begin{abstract}
Este trabalho apresenta o desenvolvimento e a avaliação científica da plataforma
AgroLeilões, um sistema web para realização de leilões agropecuários com suporte a
transmissão ao vivo e lances em tempo real. O problema investigado é a ineficiência
logística e a ausência de governança nos canais informais utilizados pelo setor, como
grupos de aplicativos de mensagens, que impõem latência operacional crítica durante os
pregões. Fundamentado na metodologia \textit{Design Science Research} (DSR), o trabalho
propõe um artefato composto por quatro perfis de usuário --- Comprador, Produtor,
Leiloeiro e Administrador --- integrados por protocolos WebSockets e WebRTC. A avaliação
foi conduzida por meio da \textit{System Usability Scale} (SUS) e da Avaliação
Heurística de Nielsen. Os resultados evidenciam usabilidade satisfatória para os perfis
Comprador (SUS 70,0--90,0) e Leiloeiro (77,5), com falha crítica no perfil Produtor
rural (SUS 55,0, abaixo do limiar de aceitabilidade de 68 pontos), correlacionada à
violação H8 de sobrecarga cognitiva no formulário de cadastro de lotes. Os achados
fundamentam princípios de design para interfaces agrotecnológicas voltadas a usuários
rurais de baixo letramento digital. O projeto enquadra-se nos preceitos da
curricularização da extensão universitária (Resolução CNE/CES
n\textsuperscript{o} 7/2018), ao disponibilizar uma solução tecnológica aplicada às
demandas reais do agronegócio brasileiro.
\end{abstract}

\begin{resumo}
Este trabalho apresenta o desenvolvimento e a avaliação científica da plataforma
AgroLeilões, um sistema web para realização de leilões agropecuários com suporte a
transmissão ao vivo e lances em tempo real. O problema investigado é a ineficiência
logística e a ausência de governança nos canais informais utilizados pelo setor, como
grupos de aplicativos de mensagens, que impõem latência operacional crítica durante os
pregões. Fundamentado na metodologia \textit{Design Science Research} (DSR), o trabalho
propõe um artefato composto por quatro perfis de usuário --- Comprador, Produtor,
Leiloeiro e Administrador --- integrados por protocolos WebSockets e WebRTC. A avaliação foi conduzida por meio da \textit{System Usability Scale} (SUS) e da
Avaliação Heurística de Nielsen. Os resultados evidenciam usabilidade satisfatória para
os perfis Comprador (SUS 70,0--90,0) e Leiloeiro (77,5), com falha crítica no perfil
Produtor rural (SUS 55,0, abaixo do limiar de aceitabilidade de 68 pontos), correlacionada
à violação H8 de sobrecarga cognitiva no formulário de cadastro de lotes. Os achados
fundamentam princípios de design para interfaces agrotecnológicas voltadas a usuários
rurais de baixo letramento digital. O projeto enquadra-se nos preceitos da
curricularização da extensão universitária (Resolução CNE/CES
n\textsuperscript{o} 7/2018), ao disponibilizar uma solução tecnológica aplicada às
demandas reais do agronegócio brasileiro.

\textbf{Palavras-chave:} leilões agropecuários, WebSockets, WebRTC, Design Science
Research, usabilidade, SUS, avaliação heurística.
\end{resumo}

% -------------------------------------------------------
% 1. INTRODUÇÃO
% -------------------------------------------------------
\section{Introdução}

O Brasil registrou exportação recorde de mais de 411,8 mil cabeças de gado vivo em
período recente, consolidando-se como maior exportador global do setor \cite{leilosul2024}.
Nesse contexto, os leilões agropecuários constituem um dos principais mecanismos de
formação de preço e liquidez do agronegócio brasileiro. Contudo, o modelo predominante
ainda exige presença física dos participantes, gerando custos logísticos expressivos ---
transporte de animais, deslocamento de compradores, aluguel de espaço físico --- além de
restringir geograficamente a base de potenciais compradores.

Diante da ausência de plataformas verticalizadas para esse nicho, parte do setor
recorreu a soluções improvisadas, como grupos de WhatsApp e transmissões informais em
redes sociais. Essas alternativas apresentam limitações críticas: ausência de controle
formal de lances, falta de rastreabilidade de transações, impossibilidade de auditoria
e latência operacional que compromete a integridade do pregão em tempo real.

Este trabalho propõe e avalia a plataforma AgroLeilões como resposta a esse problema.
A hipótese central é que uma plataforma web verticalizada, estruturada sobre protocolos
de comunicação síncrona (WebSockets e WebRTC) e com controle de acesso baseado em
perfis de usuário, é capaz de reduzir a latência transacional, aumentar a
rastreabilidade dos leilões e expandir geograficamente a participação de compradores e
produtores rurais.

Do ponto de vista extensionista, o projeto materializa os preceitos da Resolução
CNE/CES n\textsuperscript{o} 7/2018 \cite{brasil2018}, que estabelece a curricularização
da extensão universitária ao exigir que projetos de graduação promovam transferência de
conhecimento técnico-científico à sociedade.

Este artigo está organizado da seguinte forma: a Seção~\ref{sec:fundamentacao} apresenta
a fundamentação teórica; a Seção~\ref{sec:metodologia} descreve os procedimentos
metodológicos; a Seção~\ref{sec:artefato} apresenta o artefato desenvolvido; a
Seção~\ref{sec:resultados} expõe os resultados e discussões; e a
Seção~\ref{sec:conclusao} apresenta as conclusões, limitações e trabalhos futuros.

% -------------------------------------------------------
% 2. FUNDAMENTAÇÃO TEÓRICA
% -------------------------------------------------------
\section{Fundamentação Teórica}
\label{sec:fundamentacao}

\subsection{Tecnologias de Comunicação em Tempo Real}

A construção de sistemas de leilão online com pregão ao vivo impõe requisitos técnicos
que as arquiteturas HTTP tradicionais não conseguem satisfazer de forma eficiente. Duas
tecnologias complementares foram adotadas no AgroLeilões.

\textbf{WebSockets} é um protocolo de comunicação \textit{full-duplex} sobre uma única
conexão TCP persistente, padronizado pela RFC 6455 \cite{fette2011}. Diferentemente do
HTTP, que exige uma nova requisição para cada atualização, o WebSocket mantém o canal
aberto entre cliente e servidor, permitindo que atualizações sejam entregues ao cliente
instantaneamente. No contexto de leilões, quando um lance é registrado, todos os
participantes conectados devem receber a atualização em milissegundos.

\textbf{WebRTC} (\textit{Web Real-Time Communication}) é uma especificação do W3C
\cite{w3c2021} que permite comunicação peer-to-peer de áudio, vídeo e dados diretamente
entre navegadores. No AgroLeilões, o WebRTC viabiliza a transmissão ao vivo do pregão
conduzido pelo Leiloeiro, permitindo que compradores acompanhem visualmente a
apresentação dos animais em tempo real.

\subsection{Metodologias de Avaliação de Usabilidade}

\textbf{Design Science Research (DSR)} é uma metodologia que tem o artefato tecnológico
como objeto e contribuição científica central. Segundo \cite{wieringa2014}, a DSR opera
sobre duas macrodinâmicas: a concepção de um artefato destinado a melhorar a realidade
dos \textit{stakeholders} e a investigação empírica sobre sua capacidade de resposta no
contexto de uso. A avaliação \textit{ex-post} --- realizada após a construção do artefato
--- é metodologicamente legítima no paradigma DSR.

\textbf{System Usability Scale (SUS)}, desenvolvida por \cite{brooke1996}, é o
instrumento padrão-ouro para avaliação expedita de usabilidade. Consiste em dez
afirmações respondidas em escala Likert de cinco pontos. O escore final é calculado pela
fórmula de Brooke: para itens ímpares, subtrai-se 1 da resposta; para itens pares,
subtrai-se a resposta de 5. A soma dos dez valores normalizados é multiplicada por 2,5,
produzindo um escore entre 0 e 100. Escores acima de 68 indicam usabilidade aceitável;
acima de 80, usabilidade excelente.

\textbf{Avaliação Heurística de Nielsen} \cite{nielsen1994} é um método de inspeção
analítica no qual avaliadores verificam a interface contra dez princípios gerais de
usabilidade. Problemas identificados são classificados em escala de severidade de 0 a 4:
0 (falso positivo), 1 (cosmético), 2 (menor), 3 (maior) e 4 (catástrofe).

\textbf{GQM (\textit{Goal/Question/Metric})} \cite{basili1984} é um \textit{framework}
que estrutura a avaliação em três níveis: Meta, Questões e Métricas. Foi utilizado para
alinhar os instrumentos de avaliação ao objetivo científico do trabalho.

\subsection{Contexto de Mercado e Soluções Existentes}

O mercado brasileiro de leilões agropecuários online é atualmente atendido por plataformas
como Megaleilões e MF Leilões Rurais \cite{mfleiloes2024}, que processam volumes expressivos
de transações de bovinos de corte, leite, ovinos e equinos. Entretanto, essas soluções
consolidadas frequentemente se apoiam em arquiteturas de transmissão de vídeo de alta
latência --- tipicamente \textit{streams} RTMP embutidos via plataformas terceiras (e.g.,
YouTube Live) --- com interfaces de lances separadas da transmissão ao vivo, impondo
ao usuário a gestão cognitiva de múltiplas janelas e fluxos de informação paralelos.

A AgroLeilões se diferencia ao propor uma arquitetura integrada que une a transmissão
WebRTC e o motor de lances WebSockets em uma única interface coesa e de baixa carga
cognitiva. Enquanto as plataformas legadas fragmentam a experiência entre o canal de vídeo
e o canal de lances, a arquitetura proposta consolida ambos os fluxos em um ambiente
unificado, reduzindo a latência de ponta a ponta e eliminando a sobrecarga de alternância
de contexto (\textit{context switching}) que compromete a agilidade decisória em pregões
de alta pressão temporal. Esta diferenciação arquitetural fundamenta a relevância científica
e de mercado do artefato proposto.

% -------------------------------------------------------
% 3. PROCEDIMENTOS METODOLÓGICOS
% -------------------------------------------------------
\section{Procedimentos Metodológicos}
\label{sec:metodologia}

Esta pesquisa adota a \textbf{Design Science Research (DSR)} como metodologia central
\cite{hevner2004,peffers2007}, classificando-se como pesquisa aplicada de natureza mista
(quantitativa e qualitativa). O protótipo AgroLeilões foi concebido a partir da
observação de uma lacuna de mercado — ausência de plataformas abertas e acessíveis para
pequenos produtores rurais. No ciclo DSR formalizado neste artigo, o desenvolvimento
constitui uma \textit{Iteração Zero} exploratória; o presente trabalho formaliza a fase
de avaliação \textit{ex-post}, reconhecida como estratégia metodologicamente legítima por
\cite{venable2012}, com o objetivo de extrair \textbf{diretrizes de design domínio-específicas}
para interfaces agrotecnológicas voltadas a usuários rurais de baixo letramento digital.
A contribuição científica não reside exclusivamente no artefato em si, mas no
conhecimento produzido a partir de sua avaliação com usuários representativos do domínio.

Segundo o \textit{Framework for Evaluation in Design Science Research} (FEDS)
\cite{venable2012}, a estratégia de avaliação empregada neste trabalho classifica-se
como \textbf{formativa e artificial}: \textit{formativa} porque seu propósito é
identificar falhas de design para orientar iterações futuras do artefato --- não provar
utilidade final em produção; \textit{artificial} porque testa um protótipo com dados
simulados em ambiente controlado, separado das pressões sistêmicas reais de um pregão ao
vivo. Essa classificação explícita é metodologicamente honesta e alinhada ao objetivo de
mitigação antecipada de riscos de design.

\textbf{Ética em Pesquisa.} Esta pesquisa envolveu participantes humanos exclusivamente
para avaliação de usabilidade de interface de \textit{software}, sem coleta de dados
sensíveis, identificação de sujeitos ou qualquer intervenção clínica ou psicológica. A
metodologia é formalmente isenta de registro no sistema CEP/CONEP, em conformidade com
a Resolução CNS n\textsuperscript{o} 510, de 7 de abril de 2016 \cite{cns2016}, Art. 1,
Parágrafo Único, Incisos VII e VIII, que exclui expressamente pesquisas de avaliação de
sistemas cujos dados não permitem identificação dos participantes.

A estrutura GQM que orienta a avaliação é:

\begin{itemize}
  \item \textbf{Meta:} Validar a eficácia interativa e a adequação de usabilidade da
    plataforma AgroLeilões em cenário de alta pressão transacional (leilões síncronos
    ao vivo).
  \item \textbf{Questão:} A segregação de funções em quatro perfis, acoplada à
    sincronização via WebSockets e WebRTC, oferece experiência de uso com baixa fricção
    para usuários do setor agropecuário?
  \item \textbf{Métricas:} (1) Escore global SUS; (2) quantidade e severidade de
    violações heurísticas identificadas.
\end{itemize}

\subsection{Avaliação Quantitativa --- SUS}

Foram recrutados 5 participantes voluntários externos ao ambiente acadêmico, com
perfis aproximados aos usuários-alvo da plataforma. Cada participante foi atribuído a
um dos perfis nativos do sistema: \textbf{Produtor} (cadastro de três lotes de animais),
\textbf{Leiloeiro} (configuração e ativação de transmissão ao vivo simulada) e
\textbf{Comprador} (acompanhamento de pregão e registro de lances em tempo real).
Os participantes foram caracterizados conforme a Tabela~\ref{tab:participantes}.

\begin{table}[h]
\centering
\caption{Perfil dos Participantes das Sessões SUS}
\label{tab:participantes}
\begin{tabularx}{\columnwidth}{llXl}
\toprule
\textbf{ID} & \textbf{Perfil} & \textbf{Descrição} & \textbf{Faixa etária} \\
\midrule
P1 & Comprador  & Universitário, familiarizado com e-commerce & 23 anos \\
P2 & Comprador  & Letramento digital médio, usa aplicativos bancários & 35 anos \\
P3 & Produtor   & Produtor rural, baixo letramento digital & 48 anos \\
P4 & Leiloeiro  & Familiaridade com streaming e gestão de eventos & 29 anos \\
P5 & Comprador  & Alto letramento digital, uso intenso de tecnologia & 21 anos \\
\bottomrule
\end{tabularx}
\end{table}

As sessões foram conduzidas remotamente via videochamada, com duração de
aproximadamente 15 minutos cada. Após a interação, cada participante respondeu ao
formulário SUS completo via Google Forms. Os escores foram calculados aplicando a
fórmula de Brooke. Dado o tamanho reduzido da amostra (\textit{n}=5, \textit{n}=1 por
perfil), os resultados SUS são tratados como \textbf{indicadores qualitativos e
formativos} --- e não como métricas quantitativas summativas com validade estatística
generalizável \cite{sauro2011} --- cuja função primária é guiar a triangulação com a
avaliação heurística e orientar decisões de design nas próximas iterações do artefato.

\subsection{Avaliação Qualitativa --- Heurísticas de Nielsen}

A avaliação heurística foi conduzida pelo autor, inspecionando sistematicamente as
telas principais da plataforma contra as dez heurísticas de Nielsen. As telas avaliadas
incluíram: painel de lances ao vivo (Comprador), transmissão e controle do pregão
(Leiloeiro), cadastro de lotes (Produtor) e painel administrativo (Administrador). Cada
violação foi registrada com: heurística violada, descrição, contexto e grau de
severidade (0--4).

% -------------------------------------------------------
% 4. O ARTEFATO AGROLEILÕES
% -------------------------------------------------------
\section{O Artefato AgroLeilões}
\label{sec:artefato}

\subsection{Visão Geral}

A AgroLeilões é uma plataforma web que reúne em um único ambiente: descoberta e
divulgação de leilões agropecuários, visualização de lotes e animais, acompanhamento de
pregões ao vivo, registro de lances, gestão financeira (carteira e pagamentos) e
administração da plataforma.

\subsection{Requisitos Funcionais e Não Funcionais}

A Tabela~\ref{tab:rf} apresenta os principais requisitos funcionais levantados durante
o desenvolvimento da plataforma.

\begin{table}[h]
\centering
\caption{Requisitos Funcionais da Plataforma AgroLeilões}
\label{tab:rf}
\begin{tabularx}{\columnwidth}{lX}
\toprule
\textbf{ID} & \textbf{Requisito} \\
\midrule
RF01 & Cadastro e autenticação de usuários por perfil (Comprador, Produtor, Leiloeiro, Administrador) \\
RF02 & Cadastro de leilões com título, categoria, tipo, data, localização, lotes e taxas \\
RF03 & Cadastro de lotes com dados do animal (raça, sexo, idade, peso, registro, genealogia) \\
RF04 & Transmissão ao vivo do pregão via WebRTC \\
RF05 & Registro e exibição de lances em tempo real via WebSockets \\
RF06 & Lance automático configurável por valor máximo \\
RF07 & Dashboard personalizado por perfil com histórico de atividades \\
RF08 & Carteira digital com saldo, depósitos e histórico de movimentações \\
RF09 & Sistema de favoritos para lotes e leilões \\
RF10 & Notificações em tempo real (lance superado, início/fim de leilão) \\
RF11 & Filtros avançados no catálogo de lotes (raça, preço, categoria) \\
RF12 & Painel administrativo com gestão de usuários, leilões, pagamentos e auditoria \\
\bottomrule
\end{tabularx}
\end{table}

A Tabela~\ref{tab:rnf} apresenta os requisitos não funcionais que orientaram as
decisões de arquitetura.

\begin{table}[h]
\centering
\caption{Requisitos Não Funcionais da Plataforma AgroLeilões}
\label{tab:rnf}
\begin{tabularx}{\columnwidth}{lX}
\toprule
\textbf{ID} & \textbf{Requisito} \\
\midrule
RNF01 & Latência de lances inferior a 500ms sob carga normal \\
RNF02 & Disponibilidade mínima de 99\% durante janelas de leilão ativo \\
RNF03 & Interface responsiva compatível com desktop e dispositivos móveis \\
RNF04 & Controle de acesso por perfil com autenticação segura \\
RNF05 & Rastreabilidade completa de ações sensíveis via módulo de auditoria \\
RNF06 & Transmissão de vídeo ao vivo com latência tolerável em conexões de banda moderada \\
\bottomrule
\end{tabularx}
\end{table}

\subsection{Arquitetura da Solução}

A Figura~\ref{fig:arquitetura} apresenta a arquitetura alvo de produção da AgroLeilões,
composta por três camadas: cliente, servidor e dados. O protótipo atual implementa
integralmente a camada cliente (\textit{frontend} React/Next.js 16 com TypeScript), com
estado de lances e sessões simulado via \textit{store} Zustand, validando todos os
fluxos de interface e interação por perfil. As camadas de servidor (API REST, WebSocket
Server, Signaling Server) e persistência (banco de dados relacional, armazenamento de
mídia) constituem a arquitetura de produção a ser implementada nas próximas iterações
do ciclo DSR.

\begin{figure}[h]
\centering
\begin{lstlisting}
+-----------------------------------------------+
|  CLIENTE (Browser) -- IMPLEMENTADO            |
|  React/Next.js Frontend                       |
|  +-- Zustand Store  (estado simulado, mock)   |
|  +-- WebSocket Client  [arquitetura alvo]     |
|  +-- WebRTC Client     [arquitetura alvo]     |
+-------------------+---------------------------+
                    | [arquitetura alvo]
+-------------------v---------------------------+
|  SERVIDOR -- ARQUITETURA ALVO                 |
|  +-- API REST     (autenticacao, dados, CRUD) |
|  +-- WebSocket Server  (engine de lances)     |
|  +-- Signaling Server  (WebRTC handshake)     |
+-------------------+---------------------------+
                    |
+-------------------v---------------------------+
|  DADOS -- ARQUITETURA ALVO                    |
|  +-- Banco de Dados Relacional                |
|     (usuarios, leiloes, lotes, lances)        |
|  +-- Armazenamento de midia (imagens de lotes)|
+-----------------------------------------------+
\end{lstlisting}
\caption{Arquitetura Alvo da Plataforma AgroLeilões (camada cliente implementada no protótipo)}
\label{fig:arquitetura}
\end{figure}

O fluxo de um lance em tempo real, na arquitetura alvo, ocorre da seguinte forma: o
Comprador submete um lance pelo \textit{frontend} → o WebSocket Client envia a mensagem
ao WebSocket Server → o servidor valida o lance (valor acima do atual, usuário
autenticado, lote ativo) → persiste no banco de dados → transmite a atualização para
todos os clientes conectados naquele leilão em menos de 500ms. No protótipo, esse fluxo
é simulado via \textit{store} Zustand, replicando fielmente a experiência de interface
sem dependência de infraestrutura de \textit{backend} implantada.

\subsection{Perfis de Usuário}

A Tabela~\ref{tab:perfis} apresenta os quatro perfis de usuário da plataforma e suas
respectivas responsabilidades.

\begin{table}[h]
\centering
\caption{Perfis de Usuário da Plataforma AgroLeilões}
\label{tab:perfis}
\begin{tabularx}{\columnwidth}{lX}
\toprule
\textbf{Perfil} & \textbf{Responsabilidades} \\
\midrule
Comprador     & Participar de leilões, dar lances, gerenciar carteira e pagamentos \\
Produtor      & Cadastrar animais e lotes, acompanhar vendas e faturamento \\
Leiloeiro     & Criar e gerenciar leilões, conduzir transmissão ao vivo \\
Administrador & Gerenciar plataforma, usuários, relatórios e auditoria \\
\bottomrule
\end{tabularx}
\end{table}

\subsection{Desafios Técnicos}

O desenvolvimento da AgroLeilões impôs desafios técnicos que evidenciam a complexidade
do projeto além de implementações CRUD convencionais.

\textbf{Sincronização de lances concorrentes.}
O cenário mais crítico de um leilão ao vivo é o registro simultâneo de lances por
múltiplos compradores no mesmo instante. Sem controle de concorrência, dois compradores
poderiam registrar lances com o mesmo valor ou fora de ordem temporal. A solução adotada
foi o processamento serializado de lances no WebSocket Server, com validação de estado
do lote antes de cada persistência.

\textbf{Latência WebRTC em conexões de banda variável.}
Produtores e leiloeiros em zonas rurais frequentemente operam com conexões de qualidade
variável. Foi necessário configurar parâmetros de qualidade adaptativa de vídeo e
implementar \textit{fallbacks} para cenários de degradação de sinal.

\textbf{Controle de estado do lote ativo.}
Um leilão ao vivo percorre sequencialmente múltiplos lotes. O sistema precisa garantir
que todos os clientes conectados vejam exatamente o mesmo lote ativo ao mesmo tempo.
O controle de estado foi implementado via eventos WebSocket de transição, com confirmação
de recebimento pelo cliente antes de aceitar lances no novo lote.

\textbf{Isolamento de sessões de leilão.}
A plataforma precisa suportar múltiplos leilões simultâneos sem que mensagens de um
pregão contaminem os canais de outro. A implementação de \textit{rooms} isoladas no
WebSocket Server resolveu esse problema.

% -------------------------------------------------------
% 5. RESULTADOS E DISCUSSÃO
% -------------------------------------------------------
\section{Resultados e Discussão}
\label{sec:resultados}

\subsection{System Usability Scale (SUS)}

Foram coletadas respostas de 5 participantes. A Tabela~\ref{tab:sus} apresenta os
escores individuais calculados pela fórmula de Brooke.

\begin{table}[h]
\centering
\caption{Escores Individuais --- System Usability Scale}
\label{tab:sus}
\begin{tabular}{llc}
\toprule
\textbf{Participante} & \textbf{Perfil Testado} & \textbf{Escore SUS} \\
\midrule
P1 & Comprador  & 80,0 \\
P2 & Comprador  & 70,0 \\
P3 & Produtor   & 55,0 \\
P4 & Leiloeiro  & 77,5 \\
P5 & Comprador  & 90,0 \\
\midrule
\textbf{Média} & --- & \textbf{74,5} \\
\bottomrule
\end{tabular}
\end{table}

A análise por perfil revela variância significativa entre os grupos: Compradores
obtiveram escores entre 70,0 e 90,0 pontos; Leiloeiro atingiu 77,5 pontos; e o Produtor
rural registrou 55,0 pontos --- único resultado abaixo do limiar de aceitabilidade de
68 pontos estabelecido por \cite{brooke1996}. O escore médio geral de 74,5 pontos
classifica-se como ``Boa'' na escala de Brooke; contudo, a agregação de perfis com
demandas cognitivas heterogêneas nesse único valor mascara a falha crítica no módulo
de cadastro de lotes e deve ser interpretada com cautela \cite{sauro2011}. O achado
central da avaliação quantitativa é a disparidade entre o perfil de maior letramento
digital (Comprador P5, 90,0 pontos) e o perfil de menor letramento (Produtor, 55,0
pontos), diferença de 35 pontos que evidencia que o design atual serve adequadamente
os compradores mas falha no atendimento ao segmento rural vendedor.

\subsection{Avaliação Heurística de Nielsen}

A inspeção sistemática identificou 4 violações heurísticas. A
Tabela~\ref{tab:heuristica} apresenta os problemas encontrados, classificados por
severidade.

\begin{table}[h]
\centering
\footnotesize
\caption{Violações Heurísticas Identificadas}
\label{tab:heuristica}
\begin{tabular}{p{0.4cm} p{3.0cm} p{6.5cm} p{2.4cm} p{0.6cm}}
\toprule
\textbf{\#} & \textbf{Heurística} & \textbf{Descrição} & \textbf{Tela} & \textbf{Sev.} \\
\midrule
1 & H1 --- Visibilidade do Status &
  Latência WebRTC não exibida ao Comprador durante o pregão; usuário não sabe se há
  atraso entre vídeo e lances &
  Leilão ao vivo & 2 \\[4pt]
2 & H5 --- Prevenção de Erros &
  Ausência de confirmação secundária para lances de alto valor; clique acidental gera
  compromisso de compra irreversível &
  Painel de lances & 3 \\[4pt]
3 & H8 --- Estética e Design Minimalista &
  Formulário de cadastro de lotes exibe todos os campos opcionais e obrigatórios
  sem distinção visual clara, gerando sobrecarga cognitiva para usuários
  de baixo letramento digital &
  Cadastro de lotes (Produtor) & 3 \\[4pt]
4 & H9 --- Ajuda ao Usuário para Reconhecer Erros &
  Mensagens de erro não indicam o campo específico com problema;
  exibem apenas ``Erro ao salvar'' sem orientação de correção &
  Formulários gerais & 2 \\
\bottomrule
\end{tabular}
\end{table}

Nenhuma violação de severidade 4 (catástrofe) foi identificada, indicando que o sistema
não possui falhas que impeçam completamente a execução das tarefas primárias. Entretanto,
as violações de severidade 3 (H5 e H8), segundo a taxonomia de \cite{nielsen1994b},
classificam-se como \textit{problemas maiores} que devem ser corrigidos antes de qualquer
implantação em produção \cite{nielsen1994b}. A ausência de confirmação de lances
financeiros irreversíveis (H5) e a sobrecarga cognitiva no formulário de cadastro (H8)
representam riscos operacionais que, no estado atual, inviabilizam o uso seguro da
plataforma pelo perfil Produtor.

\subsection{Discussão}

Os resultados da avaliação permitem qualificar a hipótese central do trabalho. O escore
médio SUS de 74,5 pontos confirma que a plataforma AgroLeilões oferece usabilidade
satisfatória para o conjunto de perfis avaliados, sustentando parcialmente a hipótese
de que uma solução verticalizada com WebSockets e WebRTC reduz a fricção operacional
em relação a canais informais. Os perfis Comprador apresentaram os maiores escores
(70,0 a 90,0 pontos), refletindo a objetividade do fluxo de lances e do acompanhamento
ao vivo --- funcionalidades cujo design prioriza clareza e velocidade de interação.
O perfil Leiloeiro obteve 77,5 pontos, indicando usabilidade satisfatória para usuários
com experiência prévia em transmissão digital.

O achado de maior relevância científica foi o escore do perfil Produtor (55,0 pontos),
único resultado abaixo do limiar de aceitabilidade (68 pontos). Este resultado evidencia
que o fluxo de cadastro de animais e lotes --- com múltiplos campos técnicos simultâneos
sem distinção visual entre obrigatórios e opcionais --- representa a principal barreira
de usabilidade da plataforma para o segmento rural de menor letramento digital. A
violação H8 (severidade 3) identificada na avaliação heurística corrobora este achado,
confirmando que a sobrecarga cognitiva no formulário de cadastro é o problema
prioritário a ser endereçado nas próximas iterações de desenvolvimento.

A ausência de violações de severidade 4 indica que o sistema não impede a execução
das tarefas primárias. As violações de severidade 3 (H5 e H8) concentram-se nos fluxos
de maior impacto financeiro e cognitivo, requerendo reestruturação obrigatória antes de
qualquer implantação com usuários reais.

\textbf{Diretrizes de Design Domínio-Específicas para Interfaces Agrotecnológicas.}
A triangulação dos resultados SUS e heurísticos permite a \textit{validação empírica},
no contexto agrotecnológico, de padrões de interação já estabelecidos na literatura de
IHC \cite{tidwell2011,gregor2013}: \textit{(i)} a \textit{Progressive Disclosure} --- padrão
que consiste em exibir inicialmente apenas campos obrigatórios e revelar opções avançadas
sob demanda \cite{tidwell2011} --- é um \textbf{requisito não-negociável} de inclusão
digital para o perfil Produtor rural, dado seu baixo letramento digital e a sobrecarga
cognitiva documentada (H8, severidade 3); \textit{(ii)} ações financeiras irreversíveis
devem obrigatoriamente exigir confirmação explícita em etapa separada, conforme o
princípio de prevenção de erros (H5); \textit{(iii)} mensagens de erro devem identificar
o campo específico e orientar a correção com linguagem não técnica (H9). A contribuição
científica deste trabalho não é a proposição de novos paradigmas de IHC, mas a
\textbf{validação empírica} de que esses padrões estabelecidos são requisitos funcionais
basilares --- não refinamentos opcionais de UX --- para a adoção digital sustentável no
segmento de produtores rurais de baixo letramento, em conformidade com a classificação
de conhecimento DSR de nível situado (\textit{situated implementation}) proposta por
\cite{gregor2013}.

% -------------------------------------------------------
% 6. CONCLUSÃO
% -------------------------------------------------------
\section{Conclusão}
\label{sec:conclusao}

Este trabalho apresentou o desenvolvimento e a avaliação científica da plataforma
AgroLeilões, proposta como solução à ineficiência dos canais informais de leilão
agropecuário. Fundamentado na metodologia DSR, o artefato foi avaliado por instrumentos
formais de usabilidade --- SUS e Avaliação Heurística de Nielsen --- produzindo
evidências quantitativas e qualitativas sobre sua adequação ao contexto de uso.

O achado central da avaliação é a falha crítica de usabilidade no perfil Produtor rural
(SUS 55,0 pontos, abaixo do limiar de aceitabilidade de 68), correlacionada à violação
H8 de severidade 3 no formulário de cadastro de lotes. Este resultado evidencia que o
design atual serve adequadamente os compradores e leiloeiros, mas não atende ao segmento
de menor letramento digital --- precisamente o público mais vulnerável e representativo
do agronegócio familiar brasileiro. A avaliação heurística identificou 4 violações, duas
de severidade 3 (H5 e H8), que constituem \textit{problemas maiores} a serem
obrigatoriamente corrigidos antes de qualquer implantação em produção \cite{nielsen1994b}.
O escore médio SUS de 74,5 pontos (\textit{n}=5, classificação ``Boa'' \cite{brooke1996})
deve ser interpretado em conjunto com a distribuição por perfil, não isoladamente, dado
que a média agrega grupos de demandas cognitivas heterogêneas \cite{sauro2011}.

Do ponto de vista extensionista, o cumprimento da Resolução CNE/CES
n\textsuperscript{o} 7/2018 \cite{brasil2018} transcende a mera disponibilização de
código-fonte em repositório aberto. O engajamento ativo de \textit{stakeholders}
comunitários --- produtores rurais reais e compradores externos ao ambiente acadêmico
--- nas sessões de co-avaliação do artefato constitui uma forma concreta de inserção
social, troca recíproca de saberes e transferência de tecnologia social \cite{brasil2018}.
Ao integrar o segmento rural diretamente no processo formativo de design de um artefato
destinado a democratizar o acesso a mercados agropecuários, o projeto atua como vetor
ativo de inclusão digital, satisfazendo o requisito de interação transformadora com a
sociedade previsto na regulação nacional de curricularização da extensão.

\subsection{Limitações Atuais}

\begin{itemize}
  \item Ausência de aplicativo mobile nativo; acesso via navegador com responsividade
    limitada em alguns fluxos
  \item Sem integração de \textit{gateway} de pagamento automatizado
  \item Amostra de usuários nos testes de usabilidade (\textit{n}=5) pequena, com apenas
    um participante por perfil; resultados devem ser interpretados com cautela quanto à
    generalização, pois \textit{n}=1 por perfil não controla variância individual
  \item \textbf{Validade ecológica limitada:} a avaliação foi conduzida em protótipo com
    dados simulados, sem persistência real de \textit{backend}. Os resultados medem
    usabilidade estrutural e navegacional da interface, não o comportamento do usuário sob
    as condições reais de alta pressão emocional, latência de rede e risco financeiro de
    um pregão ao vivo. Avaliações futuras com sistema em produção são necessárias para
    validar a validade externa dos achados
  \item \textbf{Instrumentos de carga cognitiva não aplicados:} dado o alto estresse
    temporal e financeiro inerente a pregões ao vivo, instrumentos como o NASA Task Load
    Index (NASA-TLX) --- que mensura demanda mental, temporal, esforço e frustração --- e
    o UMUX-Lite --- versão condensada correlacionada ao SUS que mede utilidade e
    facilidade de uso --- seriam complementos metodológicos valiosos para futuras
    iterações de avaliação em condições de campo real
\end{itemize}

\subsection{Trabalhos Futuros}

\begin{itemize}
  \item Desenvolvimento de aplicativo mobile nativo (iOS e Android)
  \item Integração com \textit{gateways} de pagamento (PIX automático, boleto)
  \item Testes de usabilidade com produtores e compradores rurais reais em condições de
    campo
  \item Relatórios analíticos de desempenho para Produtores e Leiloeiros
  \item Avaliação de desempenho sob carga (\textit{stress testing} com múltiplos
    leilões simultâneos)
\end{itemize}

% -------------------------------------------------------
% DECLARAÇÃO DE USO DE IA (CNPq Portaria 2.664/2026)
% -------------------------------------------------------
\section*{Declaração de Uso de Inteligência Artificial Generativa}

Em conformidade com a Portaria CNPq n\textsuperscript{o} 2.664/2026, declaro que as
seguintes ferramentas de Inteligência Artificial Generativa foram utilizadas no
desenvolvimento deste trabalho:

\begin{enumerate}
  \item \textbf{Ferramenta:} Claude (Anthropic) \\
        \textbf{Etapa:} Estruturação do artigo e revisão de conteúdo \\
        \textbf{Finalidade:} Sugestão de estrutura de seções e revisão de coerência
        textual. Todo o conteúdo gerado foi revisado, adaptado e validado pelo autor.
\end{enumerate}

Declaro que todo o conteúdo gerado pelas ferramentas acima foi submetido a revisão
crítica e adaptação autoral. Assumo integral responsabilidade pela originalidade,
veracidade e integridade científica deste trabalho.

% -------------------------------------------------------
% REFERÊNCIAS
% -------------------------------------------------------
% \bibliographystyle{sbc}
% \bibliography{referencias}

% Referências inline (sem arquivo .bib externo):
\begin{thebibliography}{10}

\bibitem{basili1984}
BASILI, V. R.; WEISS, D. M.
A methodology for collecting valid software engineering data.
\textbf{IEEE Transactions on Software Engineering}, v. SE-10, n. 6,
p. 728--738, 1984.

\bibitem{brooke1996}
BROOKE, J.
SUS: A quick and dirty usability scale.
In: JORDAN, P. W. et al. (Ed.).
\textbf{Usability Evaluation in Industry}.
Londres: Taylor \& Francis, 1996. p. 189--194.

\bibitem{brasil2018}
BRASIL. Ministério da Educação.
\textbf{Resolução CNE/CES n\textsuperscript{o} 7, de 18 de dezembro de 2018}.
Estabelece as Diretrizes para a Extensão na Educação Superior Brasileira.
Brasília: MEC, 2018.

\bibitem{fette2011}
FETTE, I.; MELNIKOV, A.
\textbf{The WebSocket Protocol}.
RFC 6455. Internet Engineering Task Force, 2011.
Disponível em: \url{https://tools.ietf.org/html/rfc6455}.

\bibitem{nielsen1994}
NIELSEN, J.
\textbf{Usability Engineering}.
San Francisco: Morgan Kaufmann, 1994.

\bibitem{nielsen1994b}
NIELSEN, J.
Severity ratings for usability problems.
\textbf{Nielsen Norman Group}, 1994.
Disponível em: \url{https://www.nngroup.com/articles/how-to-rate-the-severity-of-usability-problems/}.

\bibitem{w3c2021}
W3C.
\textbf{WebRTC 1.0: Real-Time Communication Between Browsers}.
W3C Recommendation, 2021.
Disponível em: \url{https://www.w3.org/TR/webrtc/}.

\bibitem{wieringa2014}
WIERINGA, R. J.
\textbf{Design Science Methodology for Information Systems and Software Engineering}.
Berlim: Springer, 2014.

\bibitem{hevner2004}
HEVNER, A. R. et al.
Design science in information systems research.
\textbf{MIS Quarterly}, v. 28, n. 1, p. 75--105, 2004.

\bibitem{peffers2007}
PEFFERS, K. et al.
A design science research methodology for information systems research.
\textbf{Journal of Management Information Systems}, v. 24, n. 3, p. 45--77, 2007.

\bibitem{venable2012}
VENABLE, J.; PRIES-HEJE, J.; BASKERVILLE, R.
A comprehensive framework for evaluation in design science research.
In: \textbf{International Conference on Design Science Research in Information Systems}.
Berlim: Springer, 2012. p. 423--438.

\bibitem{sauro2011}
SAURO, J.; LEWIS, J. R.
\textbf{Quantifying the User Experience: Practical Statistics for User Research}.
Waltham: Morgan Kaufmann, 2011.

\bibitem{gregor2013}
GREGOR, S.; HEVNER, A. R.
Positioning and presenting design science research for maximum impact.
\textbf{MIS Quarterly}, v. 37, n. 2, p. 337--355, 2013.

\bibitem{tidwell2011}
TIDWELL, J.
\textbf{Designing Interfaces: Patterns for Effective Interaction Design}.
2. ed. Sebastopol: O'Reilly Media, 2011.

\bibitem{cns2016}
BRASIL. Conselho Nacional de Saúde.
\textbf{Resolução CNS n\textsuperscript{o} 510, de 7 de abril de 2016}.
Dispõe sobre as normas aplicáveis a pesquisas em Ciências Humanas e Sociais.
Brasília: CNS, 2016.

\bibitem{leilosul2024}
LEILOSUL.
Brasil embarca mais de 411,8 mil cabeças e exportação de gado vivo bate recorde.
\textbf{Leilosul Notícias}, 2024.
Disponível em: \url{https://leilosul.com.br/noticias/brasil-embarca-mais-de-4118-mil-cabecas-e-exportacao-de-gado-vivo-bate-recorde}.

\bibitem{mfleiloes2024}
MF LEILÕES RURAIS.
\textbf{Plataforma de Leilões Online de Animais}.
2024.
Disponível em: \url{https://www.mfleiloes.com.br/}.

\end{thebibliography}

\end{document}
